\section{The object being studied}
\label{sec:problem}

\subsection{Three different objects}

An NCA is a dynamical system whose local transition rule is parametrized by a
neural network. A graph specifies which cells exchange information.
A transformer is a separate predictor trained on serialized trajectories.
These roles must be kept distinct. The NCA background is
\cite{mordvintsev2020}; the message-passing viewpoint is
\cite{gilmer2017}. Synthetic-then-language training follows the research
question of \cite{lee2026}, without inheriting that paper's positive
empirical conclusions.

Let $G$ have $n$ vertices, let the cell alphabet be
$\mathcal A=\{0,\ldots,q-1\}$, and write
\begin{equation}
 X^t=(X_1^t,\ldots,X_n^t)\in\mathcal A^n,\qquad
 X^{t+1}=F_{\omega,G}(X^t).
 \label{eq:dynamics}
\end{equation}
Here $\omega$ is sampled once per trajectory and then held fixed. It is not
updated by the transformer's optimizer. A serialization map $\mathcal S$
produces a token sequence $Z=\mathcal S(X^{0:T-1})$. Sampling
$\omega$, $X^0$, and, if declared, $G$, defines a synthetic distribution
$\mathcal D_G$. Graph versus grid pre-pre-training changes $\mathcal D_G$.
It need not change the transformer architecture.

The binary pilot uses deterministic argmax transitions. A stochastic
extension instead specifies a transition kernel
$p_\omega(X^{t+1}\mid X^t,G)$. Statements below that require determinism
do not automatically apply to that extension.

\subsection{The primary estimand}

Partition transformer parameters into a backbone $B$ and a token
embedding/output table $E$. With tied embeddings, the same $E$ appears at
the input and output. Position parameters belong to $B$ here.
Let $\mathcal U_{G,K}$ be $K$ synthetic optimizer updates and let
\begin{equation}
 \mathcal T(B,E)=(B,E_0)
 \label{eq:reset}
\end{equation}
reset the token table to a common language initialization $E_0$.
Writing $(B_G^{\rm ppt},E_G^{\rm ppt})
=\mathcal U_{G,K}(B_0,E_{{\rm syn},0})$, the language-stage initializations are
\begin{equation}
 \begin{aligned}
 \theta_G^0&=(B_G^{\rm ppt},E_0),\\
 \theta_{\rm grid}^0&=(B_{\rm grid}^{\rm ppt},E_0),\\
 \theta_{\rm scratch}^0&=(B_0,E_0).
 \end{aligned}
 \label{eq:initializations}
\end{equation}
The synthetic initial table $E_{{\rm syn},0}$ is common to matched synthetic
arms; its vocabulary can differ from the language table's vocabulary.
The binary pilot happens to use a common 274-symbol allocation.
All arms then receive the same language update map $\mathcal V_n$,
including optimizer initialization, learning-rate schedule and minibatch
order. A population comparison at language update $n$ is
\begin{equation}
 \Delta_G(n)=
 \E\!\left[
 L_{\rm lang}(\mathcal V_n(\theta_G^0))-
 L_{\rm lang}(\mathcal V_n(\theta_{\rm scratch}^0))
 \right].
 \label{eq:estimand}
\end{equation}
The expectation averages declared paired initializations and data/rule/graph
sampling. A negative value is favorable in loss units. The graph-specific
effect replaces scratch in \eqref{eq:estimand} by the grid treatment.
These are different questions.

The historical fixed-total-budget design instead compares $K+N$ text
updates against $K$ synthetic plus $N$ text updates, with an additional
scratch scheduler boundary. That design does not identify
\eqref{eq:estimand}. This note derives conditions for the matched
language-stage question, rather than relabeling historical measurements.

\subsection{A loss and two clocks}

For a language distribution and autoregressive model $p_\theta$, define
\begin{equation}
 L_{\rm lang}(\theta)
 =\E\!\left[-\frac{1}{m}\sum_{j=1}^{m}
       \log p_\theta(Y_j\mid Y_{<j})\right].
\end{equation}
Logs are natural unless a quantity is explicitly in bits.
If $\varepsilon$ is a predeclared loss target, define the first
language-stage hit
\begin{equation}
 N_G(\varepsilon)=\inf\{n\geq0:L_{\rm lang}(\mathcal V_n(\theta_G^0))
                               \leq\varepsilon\}.
 \label{eq:hit}
\end{equation}
An empty set means $+\infty$, not zero speedup. Actual validation is
sampled and evaluated at a finite cadence, so observed hitting times are
coarser than \eqref{eq:hit}.

Language-stage cost and total preparation-plus-training cost are separate:
\begin{align}
 C_G^{\rm lang}&=c_G^{\rm lang}N_G,\\
 C_G^{\rm total}&=C_G^{\rm gen}+c_G^{\rm ppt}K+
                         c_G^{\rm lang}N_G.
 \label{eq:twocosts}
\end{align}
The constants must use a common unit, such as seconds or token-equivalent
cost. We solve their break-even condition in Section~\ref{sec:numerical}.

\subsection{Notation}

\begin{center}
\begin{tabular}{@{}p{.20\linewidth}p{.73\linewidth}@{}}
\toprule
Symbol & Meaning\\
\midrule
$n,q,d,T$ & Vertices, states/cell, non-self relation types, retained frames.\\
$P_r$ & Permutation matrix gathering the source for relation $r$.\\
$\omega,F$ & Neural rule parameters and their induced discrete rule.\\
$B,E$ & Transformer backbone and token embedding/output parameters.\\
$K,N$ & Synthetic updates and common language-stage updates.\\
$P,\ell,W$ & Patches/frame, tokens/frame, preceding-token context length.\\
$H,\mu,L$ & Local language Hessian and its lower/upper eigenvalue bounds.\\
\bottomrule
\end{tabular}
\end{center}

